\cotefile{1}
\scanpage{79}

\reversemarginpar
\margin{\begin{tikzpicture}[x=1ex,y=1ex]
\draw[line width=.65pt] (.42,2.4) .. controls (.32,2.15) and (.33,1.95) .. (.53,1.82) -- (.57,.25) .. controls (.57,.05) and (.66,0) .. (.73,.04);
\draw[line width=.65pt] (.45,1.94) .. controls (.19,2.04) and (.12,1.91) .. (.15,1.67);
\end{tikzpicture}}
\normalmarginpar
\uline{Corollary 4} Let $u$ be a Fredholm operator on the Hilbert space $E$. Then

\[
\text{(18)}\qquad {\color{red}\boxed{\color{black}|\alpha_n(u)|\leq\alpha_n(\deleted{\unsure{\sqrt{u^*u}}}|u|\deleted{{}_{\unsure{2}}})}}
\]

(\unsure{\deleted{In}|In} other words) the coefficients of the Fredholm expansion $\det(1+zu)$ are bounded in absolute value, term by term, by the coefficients of the Fredholm determinant $\det(1+z\unsure{h|H|\lvert u\rvert})$ relative to the positive Hermitian operator $h=\sqrt{u^*u}$. We indeed conclude from this

\uline{Corollary 5} Under the preceding conditions, we have

\[
\text{(19)}\qquad |\det(1+zu)|\leq\det(1+|z|h)\quad\deleted{\unsure{(\det\sqrt{u^*u})}}
\]

\unsure{Let us show|Let us prove} (18). We have \inserted{indeed}, $|\alpha_nu|=|\mathrm{Tr}\,\overset{n}{\Lambda}u|\leq\|\overset{n}{\Lambda}u\|_1=\alpha_n(\sqrt{u^*u})$ (formula (16)).

\deleted{\uline{Theorem 2} Let $u$ be a Fredholm \deleted{map}\inserted{operator} on the Hilbert space $H$. Then the sequence \inserted{$(\lambda_i)$} of eigenvalues \deleted{\illegible{1}} of $u$ is summable \deleted{\unsure{and}} \deleted{\unsure{(1)|(19)}} \deleted{$\unsure{\sum_i|\lambda_i|\leq\lVert u\rVert_1}$} \deleted{\unsure{values of $u$}}, we have}

\[
\text{\deleted{(20)}}\qquad {\color{red}\boxed{\color{black}\deleted{\det(1+u)=\prod_i(1+\lambda_i)}}}
\]

\deleted{\inserted{Applying (21) to $zu$ instead of $u$, we find}}

\deleted{\uline{Corollary 1} The Fredholm determinant of $\det(1+zu)$ of $u$ is an entire function of genus $0$\,;}

\[
\text{\deleted{(22)}}\qquad \deleted{\det(1+zu)=\prod_i(1+z\lambda_i)}
\]

\reversemarginpar
\margin{Corollary 1 ci-\unsure{$=$} is equivalent \deleted{\unsure{to the}}\inserted{to the} formula,\par
$\text{(22)}\quad\alpha_n(u)=\sum_{i_1<\mathord{-}<i_n}\lambda_{i_1}\deleted{\unsure{(u)}}\mathord{-}\lambda_{i_n}\deleted{\unsure{(u)}}$\par
of which the most interesting is
\begin{tikzpicture}[x=1ex,y=1ex,baseline=0ex]
\draw[line width=.55pt] (.05,1.55) .. controls (.45,1.65) and (.66,1.83) .. (.91,1.99);
\draw[line width=.55pt] (1.45,1.39) -- (.58,.65) -- (.77,.52);
\end{tikzpicture}\par
$\text{(23)}\quad{\color{red}\boxed{\color{black}\mathrm{Tr}\,u=\sum\lambda_i(u)}}$}
\normalmarginpar

\deleted{\uline{Proof} of theorem 2. \deleted{\unsure{We prove}} \unsure{We first prove the lemma} (20) \unsure{on the values}, \unsure{restricting ourselves to} the formula (20) : \unsure{let a} \illegible{4}, \illegible{3} \unsure{proper}, \illegible{5}, \unsure{then it all comes down to its being so}, \unsure{and that} \illegible{1} \unsure{is identical as soon as} \deleted{\unsure{all}} \unsure{the} \illegible{1} \unsure{is zero}.}

\margin{80}
